% =====================================================================
%  chunks/ch03i.tex —— 原书 p.113–124（PDF 120–131）
%  第 3 章 §8 K 理论的定义与性质；Atiyah--Hirzebruch 谱序列；
%  Adams 运算；Thom 同构与 Riemann--Roch 定理的类似物；
%  椭圆算子与 K 理论；变换群；四维流形
% =====================================================================
% [原文 p.113]
\section[K 理论。 Atiyah--Hirzebruch 谱序列。 椭圆算子与 K 理论]{$K$ 理论的定义与性质。Atiyah--Hirzebruch 谱序列。Adams 运算。Thom 同构与 Riemann--Roch 定理的类似物。椭圆算子与 $K$ 理论。变换群。四维流形}

\takeaway{本节把上同调论的两大推广之一——$K$ 理论——从向量丛的稳定等价类讲起：先是 Grothendieck 群与 Bott 周期性给出的 $\Z_2$（复情形）与 $\Z_8$（实情形）分次，其次是用 Atiyah--Hirzebruch 谱序列由上同调算出 $K$ 群，再是 Adams 运算 $\psi^k$ 与 $J$ 同态，最后经由 Thom 同构的 $K$ 理论类似物到达 Riemann--Roch 型定理与 Atiyah--Singer 指标定理。}

上同调论的最初的、也是最重要的两个推广——$K$ 理论与配边理论——是 Atiyah 在 20 世纪 50 年代末与 60 年代引入的（更确切地说，是首先从恰当的观点加以考察的）。与之相关的方法论随后的发展，极大地扩充了拓扑学的代数工具。此外，对许多拓扑问题的研究而言，$K$ 理论或配边理论都成了最恰当的框架。广义同调（与上同调）论的一般公理系统则由 G. W. Whitehead 在 20 世纪 60 年代初整理完成。

% [Tutor Note: 稳定等价 $\nu_1\oplus\varepsilon^{N_1}=\nu_2\oplus\varepsilon^{N_2}$ 里两边的 $N_1,N_2$ 可以不同，这正是「稳定」二字的含义；Grothendieck 群的元素写成形式差 $[\nu]-[\varepsilon^N]$，而 (8.2) 说明这个形式差确有几何实现。]
我们从 $K$ 理论的基本概念讲起。对有限 CW 复形对 $(K,L)$，群 $K_\R^0(K,L)$ 与 $K_\C^0(K,L)$（分别为实与复情形）定义为以 $B=K/L$ 为底空间的\emph{稳定等价的向量丛}的等价类的 \emph{Grothendieck 群}，其中两个向量丛 $\nu_1$ 与 $\nu_2$ 的「稳定等价」是指
\begin{equation}
\nu_1\oplus\varepsilon^{N_1}=\nu_2\oplus\varepsilon^{N_2}, \tag{8.1}
\end{equation}
这里 $\varepsilon^{N_i}$（$i=1,2$）是以 $\R^{N_i}$ 或 $\C^{N_i}$ 为纤维的平凡向量丛。$K/L$ 上的每个向量丛 $\nu$ 都有稳定逆。这一点在第 4 章 §1 中解释过：先把 $K/L$ 实现为流形 $U\supset K/L$ 的形变收缩核，$\nu$ 可延拓为 $U$ 上稳定等价于 $U$ 的切丛的丛；然后，逆元 $-\nu$ 实现为 $U$ 关于嵌入 $U\subset\R^q$（$q$ 充分大）的法丛：
\begin{equation}
\nu\oplus(-\nu)=\varepsilon^N\sim0. \tag{8.2}
\end{equation}
显然，在向量丛的直和运算下，$K_\Lambda^0(K,L)$（$\Lambda=\R,\C$）成群；而经张量积，它事实上成为交换环。

对有限 CW 复形 $K$，令
\begin{equation}
K_\Lambda^0(K)=K_\Lambda^0(K\cup*,*),\qquad \Lambda=\R,\C, \tag{8.3}
\end{equation}
其中 $*$ 表示单点空间，即实际上我们考虑的是 $K$ 上的向量丛。

受悬挂同构（见 (3.6)）的启发，我们定义
\begin{equation}
K_\Lambda^{-j}(K,L)=K_\Lambda^0(\Sigma^jK,\Sigma^jL),\qquad j>0, \tag{8.4}
\end{equation}
其中 $\Sigma K,\Sigma L$ 表示 $K,L$ 的悬挂。由 $G=\lim\limits_{\longrightarrow}O_n$、$\lim\limits_{\longrightarrow}U_n$ 的万有 $G$-丛的性质，可以得到自然的同构
\begin{equation}
\begin{aligned}
K_\R^{-j}(K,L)&\cong[\Sigma^j(K/L),\,BO],\\
K_\C^{-j}(K,L)&\cong[\Sigma^j(K/L),\,BU],
\end{aligned}
\tag{8.5}
\end{equation}

% [原文 p.114]
其中 $[X,Y]$ 表示映射 $X\to Y$ 的同伦类的集合。于是 Bott 周期性（见第 4 章 §2）给出
\begin{equation}
\begin{aligned}
K_\R^{-j-8}(K,L)&\cong K_\R^{-j}(K,L),\\
K_\C^{-j-2}(K,L)&\cong K_\C^{-j}(K,L),
\end{aligned}
\tag{8.6}
\end{equation}
这使我们可以借助周期性对一切 $j$（$-\infty<j<\infty$）定义群 $K_\Lambda^j(K,L)$（$\Lambda=\R,\C$）。

函子 $K_\Lambda^j(\cdot,\cdot)$ 满足广义上同调论的公理（见 §3）：同伦不变性、函子性、切除公理（因为由定义有 $K_\Lambda^j(K/L,*)=K_\Lambda^j(K,L)$），以及对偶正合列的存在性：
\begin{equation}
\cdots\longrightarrow K^j(K,L)\longrightarrow K^j(K)\longrightarrow K^j(L)\xrightarrow{\delta}K^{j+1}(K,L)\longrightarrow\cdots . \tag{8.7}
\end{equation}
对偶的同调论 $K_j^\Lambda(\cdot,\cdot)$ 对有限复形范畴定义为
\begin{equation}
\begin{aligned}
K_j^\Lambda(K)&=K_j^\Lambda(K,*)=K_\Lambda^{n-j-1}(S^n\setminus K,*),\qquad n\to\infty,\\
K_j^\Lambda(K,L)&=K_j^\Lambda(K/L,*),
\end{aligned}
\tag{8.8}
\end{equation}
即作为适当的极限，其中有限 CW 复形 $K$ 嵌入 $S^n$（$n$ 充分大）。群 $K_j^\Lambda(K)$ 没有简单的几何描述，它们在同伦论中一般并不使用。

由于悬挂 $\Sigma K$ 上的每个向量丛都由某个映射 $K\to O_n$ 或 $K\to U_n$（$n$ 充分大）的同伦类决定，我们推得自然的同构
\[
K_\R^{-1}(K)\cong[K,O],\qquad K_\C^{-1}(K)\cong[K,U].
\]
于是对可缩的 CW 复形（特别地，对单点空间），由 Bott 周期性得到同构
\[
K_\R^{-j}(*)\cong\pi_j(BO)\cong\pi_{j-1}(O),\qquad (j>0),
\]
\[
K_\C^{-j}(*)\cong\pi_j(BU)\cong\pi_{j-1}(U),\qquad (j>0).
\]
注意一般地 $K_\Lambda^*(K,L)$ 是分次斜交换环。上述观察导致对系数环 $K_j^\R(*)$、$K_j^\C(*)$ 的下述描述：

% [原文 p.115]
\begin{enumerate}[label=\arabic*.,leftmargin=2em]
\item $K_\R^*(*)$ 的生成元为 $1\in K_\R^0(*)$、$h\in K_\R^{-1}(*)$、$u\in K_\R^{-4}(*)$、$v\in K_\R^{-8}(*)$、$v^{-1}\in K_\R^8(*)$，其关系为
\begin{equation}
2h=0,\qquad h^3=0,\qquad u^2=4v,\qquad vv^{-1}=1; \tag{8.9}
\end{equation}

\item $K_\C^*(*)$ 的生成元为
\begin{equation}
1\in K_\C^0(*),\qquad w\in K_\C^{-2}(*),\qquad w^{-1}\in K_\C^2(*), \tag{8.10}
\end{equation}
它们之间只有一个关系 $ww^{-1}=1$，从而
\begin{equation}
K_\C^*(*)\cong\Z[w,w^{-1}]. \tag{8.11}
\end{equation}
\end{enumerate}

注意在 $K_\R^*(*)$ 的情形下乘以 $v$（在 $K_\C^*(*)$ 的情形下乘以 $w$）对应于应用 Bott 周期性算子。

环 $K_\Lambda^*(K)$ 可以借助 \emph{Atiyah--Hirzebruch 谱序列}计算，这种谱序列对任何广义上同调论都存在。在目前的情形下，Atiyah--Hirzebruch 谱序列以环 $E_m^{*,*}$ 为项，其微分 $d_m$ 满足
\begin{equation}
\begin{aligned}
E_{m+1}^{*,*}&=H(E_m^{*,*},d_m),\\
d_m&:E_m^{p,q}\longrightarrow E_m^{p+m,q-m+1},\qquad d_m^2=0,\\
E_2^{p,q}&=H^p(K;K_\Lambda^q(*)),\\
\sum_{p+q=s}E_\infty^{p,q}&=GK_\Lambda^s(K),
\end{aligned}
\tag{8.12}
\end{equation}
即环 $E_\infty^{*,*}=\sum\limits_{p,q}E_\infty^{p,q}$ 是 $K_\Lambda^*(K)$ 的相伴环。微分 $d_m$ 的像都是有限群；对有限 CW 复形 $K$，当 $m$ 充分大时 $d_m=0$。群 $\Image d_m$ 的阶以及微分 $d_m$ 与上同调运算的关系，由 Buchstaber 在 20 世纪 60 年代末研究过。

% [Tutor Note: (8.12) 的 $E_2^{p,q}=H^p(K;K_\Lambda^q(*))$ 只用到系数环 $K_\Lambda^q(*)$，底空间的信息全在 $H^p$ 里；退化只保证 $E_2=E_\infty$ 作为群相加回 $K_\Lambda^s(K)$，环结构仍可能带扩张。]
对整上同调群无挠的 CW 复形 $K$，当 $\Lambda=\C$ 时 Atiyah--Hirzebruch 谱序列退化，即对 $m\ge3$ 有 $d_m=0$，从而给出同构
\[
H^*(K;K_\C^*(*))\cong K_\C^*(K).
\]
当 $\Lambda=\R$ 时，对 $K$ 的同一条件可推出一切群 $\Image d_m$ 的阶至多为 2。

\begin{example}{1}
设 $K=\RP^k$。环 $K_\R^0(\RP^k)$ 由元素 $u=\eta-1$ 生成，其中 $\eta$ 是 $\RP^k$ 上的标准线丛（其非平凡 Stiefel--Whitney 类为 $w_1(\eta)$；见第 4 章 §1）。

% [原文 p.116]
由此可知，在 Atiyah--Hirzebruch 谱序列中对 $m\ge3$ 有 $d_m=0$，于是
\begin{equation}
\begin{aligned}
GK_\R^0(\RP^k,*)&\cong\sum_{j\ge0}H^j(\RP^k;K_\R^{-j}(*)),\\
K_\R^0(\RP^k,*)&\cong\Z_{2^{h(k)}},
\end{aligned}
\tag{8.13}
\end{equation}
其中
\[
k=8q+s,\qquad h(k)=4q+\lambda_s,
\]
\[
\lambda_1=1,\quad\lambda_2=2,\quad\lambda_3=2,\quad\lambda_4=3,\quad\lambda_5=3,\quad\lambda_6=3,\quad\lambda_7=3.
\]
成立下述关系：
\[
-u^2=2u,\qquad 2^{h(k)}u=0,\qquad \psi^l(u)=(1+u)^l-1,
\]
其中 $\psi^l$ 是「Adams 运算」，其定义不久将给出。

对 $\Lambda=\C$，$K_\C^0(\RP^k)$ 的生成元是 $v=cu$，即 $u$ 的复化，于是同样在 Atiyah--Hirzebruch 谱序列中对 $m\ge3$ 有 $d_m=0$，从而给出
\begin{equation}
K_\C^0(\RP^k)\cong\Z_{2^q},\qquad -v^2=2v,\qquad \psi^l(v)=(1+v)^l-1. \tag{8.14}
\end{equation}
\end{example}

借助实表示与复表示，可以在 Grothendieck 群 $K_\Lambda^0$（$\Lambda=\R,\C$）上定义各种与连续映射交换的运算。这些运算是通常上同调运算的 $K$ 理论类似物。下面定义复 $K$ 理论的 \emph{Adams 运算}。

给定 $K/L$ 上的复丛 $\eta$，可以构造外幂丛 $\Lambda^j(\eta)$（$j\ge0$）。$t$ 的形式幂级数
\[
\Lambda_t(\eta)=\sum_{j\ge0}\Lambda^j(\eta)t^j=1+\Lambda^1(\eta)t+\Lambda^2(\eta)t^2+\cdots
\]
定义了一个算子
\[
\Lambda_t:K^0(K,L)\longrightarrow K^0(K,L)[[t]],
\]
然而其系数 $\Lambda^j$ 不是可加的，因为（见第 4 章 §1）
\begin{equation}
\Lambda_t(\eta_1\oplus\eta_2)=\Lambda_t(\eta_1)\Lambda_t(\eta_2). \tag{8.15}
\end{equation}
为得到可加的运算，考虑级数
\[
-\log(1-t\Lambda^1+\cdots)=-\log\Lambda_{-t}=\sum_{k\ge1}\frac{\psi^k}{k}t^k
\]
的系数。这里 $\psi^k(\Lambda^1,\dots,\Lambda^k)$ 是 $\Lambda^j$ 的整系数多项式。用对称多项式来表示 $\psi^k$ 比较方便。设 $s_1,s_2,\dots$ 是 $u_1,u_2,\dots$ 的初等对称多项式；则对称多项式 $u_1^k+\cdots+u_k^k$ 可表为 $s_1,\dots,s_k$ 的多项式：
\[
u_1^k+\cdots+u_k^k=P_k(s_1,s_2,\dots,s_k),\qquad k=1,2,\dots .
\]

% [原文 p.117]
下述表达式（\emph{第 $k$ 个 Adams 运算}）对具体计算非常方便：
\[
\psi^k(\eta)=P_k(\eta,\Lambda^1\eta,\dots,\Lambda^k\eta).
\]
由此可知
\begin{equation}
\psi^k(\eta_1\oplus\eta_2)=\psi^k(\eta_1)\oplus\psi^k(\eta_2),\qquad \psi^1=\Lambda^1=1. \tag{8.16}
\end{equation}
若 $\eta$ 是线丛，则
\begin{equation}
\psi^k(\eta)=\eta^k. \tag{8.17}
\end{equation}
对线丛（看作 $K_\C^0$ 的元素）以及进而对它们的直和，容易验证下述性质：
\begin{equation}
\begin{aligned}
\psi^k(\eta_1\eta_2)&=\psi^k(\eta_1)\psi^k(\eta_2),\\
ch^n(\psi^k(\eta))&=k^n(ch^n\eta),\\
\psi^k\psi^l&=\psi^{kl},
\end{aligned}
\tag{8.18}
\end{equation}
其中 $ch$ 是 Chern 特征。这些性质随后不难推广到一切向量丛（或更确切地说，$K_\C^0$ 的元素）。

用 $\eta$ 表示 $K_\C^0(S^{2n},*)\cong\Z$ 的生成元，使 $ch\,\eta=\mu\in H^{2n}(S^{2n};\Z)$ 为基本同调类的对偶，则
\begin{equation}
\psi^k\eta=k^n\eta. \tag{8.19}
\end{equation}
公式 (8.19) 蕴涵下述关于算子 $\psi^k$ 与 Bott 周期性算子 $w$（见上）的交换规则：
\[
\psi^k(wz)=kw\psi^k(z).
\]
这促使我们通过令 $\psi^kw=kw$ 并利用乘性，把 $\psi^k$ 延拓到一切 $K_\C^{-j}(K,L)$（$j>0$）上：
\[
\psi^k(ab)=\psi^k(a)\psi^k(b).
\]
把算子 $\psi^k$ 延拓到 $j>0$ 的 $K_\C^j$ 上，只有在与 $\Z\bigl[\tfrac1k\bigr]$ 作张量积之后才可能，即必须过渡到群 $K_\C^j\otimes\Z\bigl[\tfrac1k\bigr]$。

在实 $K$ 理论中，Adams 运算 $\psi^k$ 以类似方式定义。这些运算与复化同态交换：
\[
c:K_\R^0(K,L)\longrightarrow K_\C^0(K,L),\qquad \psi^k\cdot c=c\cdot\psi^k.
\]

% [原文 p.118]
Adams 在 20 世纪 60 年代初引入的运算 $\psi^k$，在若干问题的解决中被证明极为有用。为叙述这些问题，我们首先需要一些进一步的定义。

设 $S(\eta)$ 表示与向量丛 $\eta$ 相伴的球面纤维化，即 $S(\eta)$ 与 $\eta$ 有相同的底空间，但向量丛 $\eta$ 的每个纤维 $\R^n$ 被替换为球面 $S^{n-1}\subset\R^n$。若在相伴球面纤维化的全空间之间存在映射
\[
S(\eta_1)\longrightarrow S(\eta_2),
\]
它在底空间上诱导恒等映射、在纤维上次数为 $+1$，则称同一底 $K$ 上纤维为 $\R^n$ 或 $\C^n$ 的两个向量丛 $\eta_1,\eta_2$ 是\emph{逐纤维同伦等价}的。由此引出实或复向量丛的\emph{稳定逐纤维同伦等价}概念（即：若存在平凡丛 $\varepsilon^{N_1},\varepsilon^{N_2}$，使 $\eta_1\oplus\varepsilon^{N_1}$ 与 $\eta_2\oplus\varepsilon^{N_2}$ 按上述意义逐纤维同伦等价，则称 $\eta_1\sim\eta_2$）。由这种稳定等价决定的类的集合（最早由 Dold 在 20 世纪 50 年代中考虑）记为 $J_\R(K)$（或 $J_\C(K)$）。于是有显然的满射
\[
K_\R^0(K)\longrightarrow J_\R(K),\qquad K_\C^0(K)\longrightarrow J_\C(K).
\]
设 $G_q$ 是球面 $S^q$ 的同伦自等价群（即映射 $S^q\to S^q$、次数为 1 的同伦类），并令 $G=\lim\limits_{\longrightarrow}G_q$ 为群 $G_q$ 在自然包含 $G_q\subset G_{q+1}$ 下的正向极限。自然嵌入 $O\subset G$（见第 4 章 §4）诱导映射 $j:BO\to BG$，于是
\[
K_\R^0(K)\cong[K,BO],
\]
\[
J_\R(K)=j_*(K_\R^0(K))\subset[K,BG],
\]
其中（如前）$[X,Y]$ 表示映射 $X\to Y$ 的同伦类的集合。

% [Tutor Note: $J$ 群的难点在于满射 $K_\R^0(K)\to J_\R(K)$ 的核没有直接的几何描述；Adams 的 $\psi^k$ 运算给出可计算的子群 $H$，于是 $K^0/H$ 是下界估计，而 Adams 猜想（后被证明）恰断言该估计在稳定意义下精确。]
对光滑流形 $M^n$，群 $J_\R(M^n)$ 有几项重要应用：
\begin{enumerate}[label=\arabic*.,leftmargin=2em]
\item 得到球面同伦群阶的下界（$M^n=S^n$）。
\item 研究球面上的光滑结构（$M^n=S^n$）。
\item 估计球面 $S^{2n-1}$ 上线性无关向量场的个数（$M^n=\RP^n$）。
\item 对光滑闭流形的稳定法丛进行分类（见第 4 章 §3、§4）。
\end{enumerate}
这里也宜提及下述事实，它是 Atiyah 一个结果的推论：

\noindent\emph{设 $\eta^N$ 是 $M^n$ 上的向量丛。类 $[\eta]\in J_\R(M^n)$ 含有 $M$ 上的法丛（即 $\eta$ 稳定同伦等价于法丛），当且仅当闭链 $\phi([M])\in H_{N+n}(T\eta;\Z)$ 是球面的。}

% [原文 p.119]
（这里 $T\eta$ 是 Thom 空间，$\phi$ 是 Thom 同构 $\phi:H_n(M^n)\xrightarrow{\ \cong\ }H_{N+n}(T\eta)$；若一个闭链 $x\in H_*(X;\Z)$ 属于 Hurewicz 同态 $\pi_*(X)\to H_*(X;\Z)$ 的像，则称它是\emph{球面的}。）由 Browder--Novikov 定理（见第 4 章 §4），对 $n\ne4k+2$、$n\ge5$ 的单连通流形 $M^n$，上述 $\phi([M^n])$ 为球面的条件，连同关于符号差的 Hirzebruch 公式，足以保证向量丛 $\eta$ 可实现为某个（可能不同的）与 $M^n$ 同伦等价的流形 $L$ 的法丛。

考虑 $K_\R^0(K)$ 的由一切形如
\begin{equation}
p^{N_p}(\psi^k-1)\eta \tag{8.20}
\end{equation}
的元素生成的子群 $H$，其中 $\eta$ 是 $K_\R^0(K)$ 的任一元素，$p$ 取遍一切素数，$N_p$ 取遍一切充分大的整数。按 Adams 的一个结果，商群
\[
\bar J_\R(K)\cong K_\R^0(K)/\langle p^{N_p}(\psi^k-1)\eta\rangle=K_\R^0(K)/H
\]
给出 $J_\R(K)$ 的下界估计，因为满射 $K_\R^0(K)\to J_\R(K)$ 的核包含在 Adams 子群 $H$ 中。在 $K=S^n$ 的情形，Adams 得以建立反向的估计（在某些维数上差因子 2），从而得到群 $J_\R(S^{4n})\subset\pi^s_{4n-1}=\pi_{4n-1+N}(S^N)$ 的阶的度量（在所指出的精确度之内）。Adams 的这一估计与 Kervaire--Milnor 得到的下界（对偶数 $n$，见第 4 章 §3）一致（差因子 2）。对 $K=\RP^{2n}$，Adams 的下界估计给出
\[
K_\R^0(\RP^{2n})\cong J_\R(\RP^{2n}),
\]
由此借助 Toda 的某些归约，可得到奇维球面 $S^{2q-1}$ 上线性无关向量场最大个数的精确值。设 $2q=2^m(2l+1)$，$m=c+4d$，其中 $0\le c\le3$，则此个数为 $2^c+8d-1$（Adams，20 世纪 60 年代初）；而且这个数目的线性无关向量场可以借助无穷小旋转在球面上构造出来。

Adams 的猜想——对充分大的 $N$，$K_\R^0(K)$ 的一切形如 $k^N(\psi^k-1)\eta$ 的元素都是 $J$-平凡的（即落在满射 $K_\R^0(K)\to J_\R(K)$ 的核中）——于 20 世纪 60 年代末被 Sullivan 与 Quillen 用涉及同伦范畴「增广」的优美的一般范畴论想法证明，这种增广使乘以 $k$ 成为可逆的。

由于同调 $K$ 理论没有简单的几何解释，闭流形的 Poincaré 对偶在 $K$ 理论中的恰当类似物需要特殊处理。特别地，出现下述问题：$K_n(M^n)$ 中的「基本类」是什么？它何时存在？一个本质等价但更易入手的问题是：对 $K$ 上纤维为 $\R^n$ 的向量丛 $\eta$，何时有「Thom 同构」
\begin{equation}
\phi_K:K_\Lambda^j(K)\longrightarrow K_\Lambda^{n+j}(T\eta),\qquad \phi_K(z)=z\phi_K(1),\qquad \Lambda=\R,\C\ ? \tag{8.21}
\end{equation}

% [原文 p.120]
对纤维为 $\C^m$（$n=2m$）的复向量丛 $\eta$，$\phi_K(1)$ 的值（在复 $K$ 理论即 $\Lambda=\C$ 的框架中，由 Grothendieck--Atiyah--Hirzebruch 于 20 世纪 50 年代末确定）为
\[
\phi_K(1)=\Lambda^{\mathrm{even}}(\eta)-\Lambda^{\mathrm{odd}}(\eta)\in K_\C^0(T\eta),
\]
其中
\[
\Lambda^{\mathrm{even}}(\eta)=\sum_{j\ge0}\Lambda^{2j}(\eta),\qquad \Lambda^{\mathrm{odd}}(\eta)=\sum_{j\ge0}\Lambda^{2j+1}(\eta).
\]
下述公式成立：
\[
\phi_H^{-1}(ch(\phi_K(1)))=T(\eta)\in H^*(K;\Q),\qquad i:K\longrightarrow T\eta,
\]
其中 $T(\eta)$ 是 $\eta$ 的「Todd 亏格」，$\phi_H$ 是上同调中的 Thom 同构（见第 4 章 §3）。由此得到 Chern 特征的公式：
\[
ch(\phi_K(\eta))=T(\eta)ch(\eta).
\]
对拟复流形（即其稳定法丛带有复结构的流形）之间的映射 $f:M^n\to N^k$，存在 Gysin 同态的类似物，其形式为
\[
Df_*D=f_!:K^0(M^n)\longrightarrow K^0(N^k),
\]
它满足
\begin{equation}
ch(f_!(\eta))T(N^k)=f_!(ch(\eta)T(M^n)). \tag{8.22}
\end{equation}
（这是 Riemann--Roch 定理的推广。）对嵌入 $M^n\subset N^k$，算子 $f_!$ 与嵌入的法丛 $\nu_{M^n\subset N^k}$ 的 Thom 同构一致。

在实 $K$ 理论（$\Lambda=\R$）中，Thom 同构的构造要求丛 $\eta$（$\dim\eta=n$）上存在 spin 结构，即要求 Stiefel--Whitney 类 $w_1(\eta),w_2(\eta)=0$。此时定义
\begin{equation}
\phi_K(1)=\Delta^+(\eta)-\Delta^-(\eta)\in K_\R^0(T\eta), \tag{8.23}
\end{equation}
其中 $\Delta^\pm$ 是群 $Spin(n)$ 在向量空间 $\Lambda(\R^n)$ 上的「半旋」表示，$\Delta^\pm(\eta)$ 是与这些表示相伴的向量丛。实 $K$ 理论中的 Thom 同构 $\phi_K$ 用来定义 $\hat A$-亏格：
\begin{equation}
\hat A(\eta)=\phi_H^{-1}(ch(\phi_K(1))). \tag{8.24}
\end{equation}
我们将会看到（在第 4 章 §3 中），对 $Spin$-流形 $M^{4k}$，$\hat A(M^{4k})=\hat A(\nu_M)$ 是整数。（这可以由广义 Riemann--Roch 定理的实类似物推得。事实上，若在 (8.22) 中取 $N^k=*$，便得到 $T$-亏格取整数值。若对拟复流形 $M$ 有 Chern 类 $c_1(M)=0$，则 $\hat A$-亏格 $\hat A(M)$ 也将是整数。）

% [原文 p.121]
经过 20 世纪 60 年代与 70 年代初若干人的研究，对大多数重要的齐性空间 $X$，环 $K_\Lambda^*(X)$ 已被算出。其中特别令人感兴趣的是下述结果：对紧连通单连通 Lie 群 $G$，环 $K_\C^*(G)$ 由「基本表示」$\rho_1,\dots,\rho_m\in K_\C^1(G)$（$m=\rank G$）在环 $\Z[w,w^{-1}]$ 上作为外代数生成：
\[
K_\C^*(G)\cong\Lambda(\rho_1,\dots,\rho_m),\qquad \Lambda=\Z[w,w^{-1}],
\]
其中 $w$ 是 Bott 周期性算子（Hodgkin 定理，20 世纪 60 年代初证明）。对 $G=U(n)$，表示 $\rho_1,\dots,\rho_n$ 恰是 $BU(n)$ 上标准丛 $\eta$ 的外幂：$\rho_j=\Lambda^j(\eta)$。对紧（或有限）Lie 群 $G$ 的分类空间 $BG$，有
\[
K_\C^{-1}(BG)=0,\qquad K_\C^0(BG)\cong\widehat R_G=\Z+I/I^2+I^2/I^3+\cdots,
\]
即复表示环 $R_G$ 关于某素理想 $I$ 的幂的完备化，其中 $R_G/I\cong\Z$（Atiyah，20 世纪 60 年代中）。

在计算无穷维 CW 复形 $X$ 的 $K$ 理论时会出现额外的困难；特别地，可能出现无穷滤过的元素，即 $K_\Lambda^*(X)$ 中限制到任一 $n$-骨架 $X^n$ 上为零的非平凡元素（Mishchenko、Buchstaber、Hodgkin，20 世纪 60 年代末）。无穷滤过元素的出现是广义上同调论的一个常见特征，在通常的理论中不会发生。

$K$ 理论最优美的应用之一与闭流形上椭圆伪微分算子的指标公式有关。先给出必要的定义。设 $\eta_1,\eta_2$ 是闭流形 $M^n$ 上的向量丛，$\Gamma(\eta_i)$ 是丛 $\eta_i$（$i=1,2$）的光滑横截面的线性空间\footnote{这里假定这些向量空间赋有恰当的范数，从而具有 $\Gamma(\eta_i)$ 上的 Sobolev 空间结构。}。$m$ 阶一般微分算子 $L$ 定义在空间 $\Gamma(\eta_1)$ 上、取值于空间 $\Gamma(\eta_2)$：
\[
L:\Gamma(\eta_1)\longrightarrow\Gamma(\eta_2).
\]
在局部坐标下 $L$ 具有形式
\[
L=\sum a_{i_1\dots i_m}(x)\frac{\partial}{\partial x^{i_1}}\cdots\frac{\partial}{\partial x^{i_m}}+L_1,
\]
其中 $a_{i_1\dots i_m}(x)$ 是矩阵值函数，$L_1$ 是阶 $<m$ 的算子。

当 $L$ 是椭圆算子时，其核与余核都是有限维的，而这两个维数之差按定义就是算子 $L$ 的\emph{指标} $\Ind L=\dim\Ker L-\dim\operatorname{Coker}L$。「指标问题」就是要把整数 $\Ind L$ 用拓扑数据表示出来，这些拓扑数据隐含在椭圆算子 $L$ 之中。

% [原文 p.122]
对算子 $L$ 中最高阶项的下述解释，提供了解决「指标问题」的关键。设 $\pi:\tau^*(M)\to M$ 是 $M^n$ 上余切丛的投影，$\tau_0^*(M)=\tau^*(M)\setminus M$ 是 $\tau^*(M)$ 中非零向量的子集。于是 $L$ 的最高阶项决定一个同态 $\sigma_L:\pi^*\eta_1\to\pi^*\eta_2$，它在局部坐标下由下式给出：
\begin{equation}
\sigma_L(x,p)=\sum a_{i_1\dots i_m}(x)p_{i_1}\cdots p_{i_m}:\C^N\longrightarrow\C^N, \tag{8.25}
\end{equation}
其中 $p=(p_1,\dots,p_n)\in\tau^*(M)$ 是余向量。同态 $\sigma_L$ 称为算子 $L$ 的\emph{符号}。算子 $L$ 是椭圆的，当且仅当其符号 $\sigma_L$ 在 $\tau_0^*(M)$ 上是同构；用局部坐标来说，当且仅当每个映射 $\sigma_L(x,p)$ 可逆，即当 $p\ne0$ 时 $\det\sigma_L(x,p)\ne0$。这个同构 $\sigma_L$——更确切地说，它的稳定同伦类 $[\sigma_L]$——应当看作椭圆算子 $L$ 的拓扑「扭」，它把算子确定到差一个紧算子为止。

对奇异积分算子有 $m=0$。在 $M=S^1$ 的情形，任何 $m=0$ 的椭圆算子 $L$ 的符号 $\sigma_L$ 归结为圆周 $S^1=\{z\in\C\mid|z|=1\}$ 上一对矩阵值函数 $\sigma_L^+(z)=\sigma_L(z,1)$、$\sigma_L^-(z)=\sigma_L(z,-1)$。此时 $L$ 的指标的 Noether--Muskhelishvili 公式（1920–1930）如下：
\[
2\pi\Ind L=\oint\log\det\sigma_L^+(z)\,dz-\oint\log\det\sigma_L^-(z)\,dz.
\]
Atiyah--Singer 的一般构造（20 世纪 60 年代初）如下：给定全空间 $\tau^*(M)$ 上的两个丛 $\pi^*(\eta_1),\pi^*(\eta_2)$，以及它们限制到子流形 $\tau_0^*(M)$ 上的同构 $\sigma_L:\pi^*(\eta_1)\to\pi^*(\eta_2)$，则定义一个\emph{差元}
\[
d(\pi^*(\eta_1),\pi^*(\eta_2),\sigma_L)\in K^0(\tau^*(M),\tau_0^*(M)).
\]
设 $D(M)$ 与 $S(M)$ 表示由 $\tau^*(M)$ 赋以某个 Riemann 度量所得到的单位圆盘丛与单位球面丛。过渡到 Thom 复形 $T\tau^*=D(M)/S(M)$（它与复形 $\tau^*(M)/\tau_0^*(M)$ 同伦等价），就得到确定好的元素
\[
d(\pi^*(\eta_1),\pi^*(\eta_2),\sigma_L)\in K^0(T\tau^*).
\]
\emph{Atiyah--Singer 指标定理}就是下述公式：
\[
\Ind L=\langle\phi^{-1}(ch\,d)\cdot T(M^n),[M^n]\rangle,
\]
其中 $T(M)$ 是 $M$ 上切丛 $\tau$ 的复化的 Todd 亏格。

这个公式的最优美应用之一，是给出 $Spin$-流形 $M^{4k}$（即满足 $w_1(M)=w_2(M)=0$ 的流形）的 $\hat A$-亏格的解释：它是 Dirac 算子的指标，从而给出 $\hat A$-亏格的整性。

% [Tutor Note: (8.25) 的符号只取最高阶项，故只依赖余向量 $p$ 与首项系数矩阵；这正是椭圆性可以在余切丛上作整体陈述、与低阶项无关的原因。]
% [原文 p.123]
目前已知这个定理有几个不同的证明以及各种推广。Atiyah--Singer 的第一个证明如同 Riemann--Roch--Hirzebruch 定理的证明，基于配边理论。这两个结果都可推广到带边流形，只要算子 $L$ 在边界上的限制是 Noether 算子。

由此可见，与线性算子理论的联系对拓扑学大有裨益，它导致若干深刻的代数-拓扑联系的发现。

作为例子，下面给出紧 Kähler 流形 $M^n$ 的全纯自映射
\[
f:M^n\longrightarrow M^n
\]
的孤立不动点个数的 Atiyah--Bott 公式。设 $p_1,\dots,p_k$ 是这样一个映射 $f$ 的不动点，全都假定是孤立且非退化的。设 $T(f):\tau\to\tau$ 是映射 $f$ 的微分，它诱导丛映射
\[
T(f):\eta_1\longrightarrow\eta_1,\qquad \eta_1=\Lambda^{\mathrm{even}}\tau,
\]
\[
T(f):\eta_2\longrightarrow\eta_2,\qquad \eta_2=\Lambda^{\mathrm{odd}}\tau.
\]
设 $\lambda_{1q},\dots,\lambda_{nq}$（$\lambda_{jq}\ne1$）表示诱导自映射 $T(f)$ 在 $M^n$ 的每个不动点 $p_q$ 处切空间 $\tau(p_q)$ 上的特征值。于是不动点个数的 \emph{Atiyah--Bott 公式}（Lefschetz 公式的类似物）如下：
\[
\chi(f)=\sum_{q=1}^{k}\Bigl(\prod_{j=1}^{n}\frac{1}{1-\lambda_{jq}}\Bigr)=\sum_{q=1}^{k}\bigl(\operatorname{Tr}T(f)|_{\eta_1(p_q)}-\operatorname{Tr}T(f)|_{\eta_2(p_q)}\bigr),
\]
其中 $\operatorname{Tr}$ 表示迹。（不动点个数的 \emph{Lefschetz 公式}是
\[
\chi(f)=\sum_m(-1)^m\operatorname{Tr}T(f)^*|_{Hol^m(M)},
\]
其中 $Hol^m(M)$ 表示 $M^n$ 上全纯 $m$-形式的向量空间。）

上述公式有几个类似物，其中包括实类似物：$M^n$ 赋以 Riemann 度量且映射 $f$ 是等距变换；或者更一般地，在 $M^n$ 上有一个 $f$-不变的椭圆复形（或 $f$-不变的椭圆算子 $L:\Gamma(\eta_1)\to\Gamma(\eta_2)$）。在后一情形，空间 $\Ker L$、$\operatorname{Coker}L$ 以及 $\eta_1,\eta_2$ 在不动点 $p_q$ 上方的纤维都成为算子 $T(f)$ 的表示空间，而下述公式（属于 Atiyah--Bott）成立：
\[
\operatorname{Tr}T(f)|_{\Ker L}-\operatorname{Tr}T(f)|_{\operatorname{Coker}L}=\sum_q\bigl(\operatorname{Tr}T(f)|_{\eta_1(p_q)}-\operatorname{Tr}T(f)|_{\eta_2(p_q)}\bigr).
\]

% [原文 p.124]
20 世纪 70 年代初，在研究 Yang--Mills 场（向量丛上的联络）的「欧氏」四维理论的过程中，发现了「瞬子」。它们是 $\R^4$ 中 Yang--Mills 方程的特殊解，当 $|x|\to\infty$ 时趋于零，因而特别地可以延拓到球面 $S^4=\R^4\cup\infty$ 上（Belavin、Polyakov、Schwarz、Tyupkin）。

特别令人感兴趣的是 Belavin--Polyakov 的自对偶方程：
\begin{equation}
\begin{aligned}
F&=F_{ab}dx^a\wedge dx^b,\qquad *F=\pm F,\\
F_{ab}&=\partial_aA_b-\partial_bA_a+[A_a,A_b],
\end{aligned}
\tag{8.26}
\end{equation}
它是对 $S^4$ 上向量丛的联络 1-形式写出的，而瞬子实际上满足它。（当然，瞬子满足完整的二阶 Maxwell--Yang--Mills 方程组，但方程 (8.26) 是一阶的。）结果表明，使泛函
\[
-S=\int_{S^4}\operatorname{Tr}\bigl(F_{ab}F^{ab}\bigr)\sqrt{g}\,d^4x=\int_{S^4}\operatorname{Tr}\bigl(F\wedge*F\bigr)
\]
全局极小的所有函数都满足自对偶方程，且对给定的第一 Pontryagin 类
\[
p_1=\mathrm{const}\cdot\int_{S^4}\operatorname{Tr}\bigl(F\wedge*F\bigr),\qquad G=SU(2).
\]

% [Tutor Note: (8.26) 是一阶方程而 Yang--Mills 方程是二阶，自对偶解自动满足后者（因为 Bianchi 恒等式给出 $d*F=0$）；把 $\R^4$ 紧化为 $S^4$ 后 $p_1$ 与泛函 $-S$ 成正比，分析问题便化成了拓扑问题。]
利用给定流形 $M^4$ 上自对偶联络的流形的性质，Donaldson 在 20 世纪 80 年代初发现了四维闭流形的深刻拓扑性质。例如，他得以证明：对单连通流形 $M^4$，由 $H_2(M^4;\Q)$ 上的相交指数定义的二次型为正定的，仅当它在 $\Z$ 上可分解为平方和。方程 (8.26) 这个值得注意的方程，在某种意义上似乎专属于四维流形的几何。这样，非线性椭圆方程的理论对拓扑学也被证明非常有用。

这里也宜提及 $K$ 理论对 $G$-空间范畴的类似物——即 Atiyah 与 Segal 在 20 世纪 60 年代中发展起来的 $K_G$ 理论。特别令人感兴趣的是带对合的空间范畴这一特殊情形（$G=\Z/2$），其中群 $\Z/2$ 通过复「反共轭」作用在向量丛的纤维上。（「$KR$ 理论」由 Atiyah 发展，其特别重要的特殊情形「$KSC$ 理论」由 Anderson 发展。）这些 $K$ 理论使人们能够得到与不动点相联系的不变量与流形的整体不变量之间的各种关系，并澄清实 $K$ 理论 8 重周期性的代数本质。
% CHUNK ch03i DONE
